\section{Funcția haotică}

Nucleul funcției hash propuse îl constituie o transformare haotică aplicată repetat asupra stării interne.
Această transformare nu se reduce la o singură hartă, ci la suprapunerea a două sisteme dinamice distincte, alese astfel încât slăbiciunile fiecăruia să fie compensate de celălalt.
Prima este harta brutarului, responsabilă de difuzie, iar a doua este harta omului din turtă dulce, responsabilă de confuzie.
Compunerea lor produce, la fiecare iterație, atât amestecarea uniformă a stării, cât și o neliniaritate suficientă pentru a împiedica orice descriere algebrică simplă a procesului. \\

\subsection{Harta brutarului}

Harta brutarului este un sistem dinamic clasic, a cărui denumire provine dintr-o analogie sugestivă cu procesul de frământare a aluatului.
Brutarul întinde aluatul pe o direcție, dublându-i lungimea, după care îl pliază la loc peste sine.
Repetarea acestui gest amestecă rapid și ireversibil orice particulă din interiorul aluatului, dispersând-o în întreaga masă.
Formal, harta acționează asupra pătratului unitate, dilatând o coordonată cu un factor de doi și contractând-o pe cealaltă cu același factor, urmate de o pliere care readuce punctul în domeniul inițial. \\

Această dublă acțiune de întindere și pliere conferă hărții brutarului proprietatea de amestecare topologică și un exponent Lyapunov strict pozitiv.
Două puncte arbitrar de apropiate sunt separate exponențial de iterările succesive, iar orbitele acoperă uniform întregul spațiu de stări, fără regiuni preferențiale sau insule de stabilitate.
În implementare, întinderea pe coordonata $x$ este realizată printr-o deplasare a biților spre stânga, contracția pe coordonata $y$ printr-o deplasare spre dreapta, iar plierea prin complementarea biților în jumătatea superioară a domeniului.
Tocmai datorită acestor proprietăți, harta brutarului joacă în cadrul funcției rolul de motor al difuziei: ea garantează că o modificare a unui singur bit din intrare se propagă rapid asupra întregii stări interne, producând efectul de avalanșă dorit. \\

Cu toate acestea, harta brutarului prezintă o limitare esențială: ea este, în esență, o transformare liniară.
Întinderea și plierea se reduc la deplasări și complementări de biți, operații afine care, luate izolat, pot fi descrise și inversate printr-un model liniar.
O construcție criptografică bazată exclusiv pe astfel de operații ar fi vulnerabilă la criptanaliza liniară, întrucât relația dintre intrare și ieșire ar rămâne, în ciuda amestecării, predictibilă din punct de vedere algebric.
Difuzia furnizată de harta brutarului este, prin urmare, necesară, dar nu și suficientă. \\

\subsection{Harta omului din turtă dulce}

Harta omului din turtă dulce, introdusă de Devaney~\cite{devaney1989}, este un sistem dinamic neliniar definit printr-o recurență simplă în care coordonata următoare depinde de valoarea absolută a coordonatei curente.
Numele său provine din forma caracteristică a atractorului, care, vizualizat în plan, amintește de silueta unei figurine din turtă dulce.
Definitorie pentru această hartă este prezența valorii absolute, singura sursă de neliniaritate, alături de cuplarea celor două coordonate: noua valoare a uneia devine vechea valoare a celeilalte, introducând astfel o memorie și o reacție inversă în dinamica sistemului. \\

Spre deosebire de harta brutarului, harta omului din turtă dulce nu poate fi exprimată ca o transformare liniară.
Valoarea absolută rupe liniaritatea, iar cuplarea coordonatelor împletește cele două jumătăți ale stării, astfel încât ieșirea nu mai poate fi descompusă într-o sumă de contribuții independente.
În cadrul funcției, acest sistem joacă rolul confuziei: el maschează orice relație statistică simplă între intrare și ieșire și introduce complexitatea algebrică pe care harta brutarului, prin natura ei liniară, nu o poate oferi. \\

Și această hartă are însă un neajuns propriu.
Spațiul ei de stări nu este uniform haotic, ci alternează regiuni de comportament haotic cu insule de stabilitate, în care orbitele devin cvasi-periodice și rămân captive.
Un punct care nimerește într-o asemenea insulă nu mai explorează întregul spațiu, ci se învârte într-un ciclu restrâns, ceea ce, într-o funcție hash, s-ar traduce prin regiuni de difuzie slabă și prin riscul apariției punctelor fixe. \\

\subsection{Suprapunerea celor două hărți}

Cele două sisteme se completează reciproc tocmai prin slăbiciunile lor complementare.
Harta brutarului oferă o amestecare uniformă și lipsită de insule de stabilitate, dar este liniară.
Harta omului din turtă dulce oferă neliniaritatea absentă din prima, dar conține regiuni cvasi-periodice în care dinamica se blochează.
Suprapunerea lor, în care fiecare rundă aplică mai întâi întinderea și plierea brutarului, iar apoi recurența neliniară a omului din turtă dulce, reunește avantajele ambelor sisteme și anulează dezavantajele fiecăruia. \\

Aplicată prima, harta brutarului dispersează punctul pe întregul spațiu de stări și îl scoate din eventualele insule de stabilitate ale celei de a doua hărți, garantând că orbita nu rămâne captivă într-un ciclu cvasi-periodic.
Aplicată în continuare, harta omului din turtă dulce împletește neliniar coordonatele astfel decorelate, conferind procesului complexitatea algebrică pe care difuzia liniară nu o poate produce singură.
Astfel, fiecare rundă a transformării realizează simultan cele două cerințe formulate de Shannon: difuzia, prin întinderea și plierea brutarului, și confuzia, prin neliniaritatea omului din turtă dulce. \\

Această rundă compusă este aplicată în mod repetat asupra fiecărei perechi de cuvinte extrase din cele două jumătăți ale stării interne, iar iterarea ei de mai multe ori amplifică efectul de amestecare până la saturație.
Rezultatul este o transformare care, deși complet deterministă și eficient calculabilă, manifestă comportamentul haotic necesar unei funcții hash robuste, întemeindu-și securitatea pe proprietățile sistemelor dinamice, iar nu pe ipoteze din teoria numerelor.

